1) Bij Asymmetrische encryptie zijn er 2 sleutels, welke zijn dit? 
Private key en Public key
 
2) Hoe noemen we een set van deze sleutels? 
Key pair
 
3) Bij RSA worden priemgetallen gebruikt, hieruit worden 2 keys berekend. Naast 
priemgetallen wordt ook nog een 3e getal gebruikt, wat is dit getal? Hoe is dit 
berekend? 
Het getal e
Euclid's algoritme, e = d^{-1} \mod \varphi(n) \quad \text{where} \quad d = e^{-1} \mod \varphi(n) \quad \text{and} \quad

1 < e < \varphi(n) \quad \text{and} \quad \gcd(e, \varphi(n)) = 1, \quad \text{where} \quad \varphi(n) = (p-1)(q-1)
 
4) Welk type encryptie ( asymmetrisch of symmetrisch ) is sneller en waarom? Welk 
type encryptie is relatief veiliger en waarom? 
 
5) Wat is een goede manier om sleutels met elkaar af te spreken, hoe heet deze 
techniek?  
 
6) Hoe heet de output van een hash functie? 
 
7) Heeft deze output een vast formaat? 
 
8) Wat kunnen wij min of meer garanderen met hashing? ( denk aan de CIA driehoek ). 
 
9) Hoe ontstaat een collision bij hashing? 
 
10) Hoe heet een bekende en zeer snelle techniek om hashing functies om te keren? Leg 
uit hoe dit werkt? 
 
11) Wat doet een checksum? En waar gebruiken wij deze voor? 
